% ===================================================================
%  GAMBAR No. 3 — Bocahan lan Nom-noman.
%
%  Traduction, faite en 2026, en javanais d'aujourd'hui, sur l'ido de
%  texto/io. Regles de la colonne : voir 10-gambar-01.tex.
%
%  DEUX SCENES, DEUX NUMEROTATIONS. Comme les tableaux 4, 8 et 9, ce
%  tableau porte deux scenes qui repartent chacune de 1, et les cles
%  le disent : t03-c1-* pour le bapteme au village, t03-c2-* pour la
%  fete publique.
%
%  LE « julio » DU BLOC t03-c2-01-1. Le fac-simile ido compose
%  « (julio od \VUgras{agosto}) » en laissant le seul juillet hors du
%  gras, alors que les deux mois qui l'encadrent y sont. C'est une
%  coquille du typographe, et gravuri/korekti.json la repare pour
%  chaque colonne. Le triple javanais — « (Juli utawa » —, a ete pose
%  AVANT que ce tableau soit ecrit : la correction attend le texte,
%  et non l'inverse. Le libelle du fichier doit donc s'y accorder a
%  la lettre.
%
%  JAVA N'A PAS DE PRINTEMPS, ET LA COLONNE LE DIT QUAND MEME. Le
%  javanais possede son propre calendrier des saisons, le pranata
%  mangsa, qui partage l'annee en douze mangsa reglees sur les pluies
%  et les travaux des champs, et il n'y a dans ce systeme ni
%  printemps ni ete. « mangsa semi » et « mangsa panas » sont les
%  termes que l'ecole javanaise emploie pour enseigner les quatre
%  saisons d'Europe : ce sont des mots de traduction, et ils sont les
%  bons, parce que la planche montre des saisons d'Europe. On
%  modernise la langue, jamais les choses — et on ne remplace pas non
%  plus la chose d'Europe par la chose d'ici.
%
%  MAIS LES JOURS ET LES MOIS, EUX, SONT JAVANAIS. Senen, Selasa,
%  Rebo — et non Senin, Selasa, Rabu comme en indonesien. Le mois se
%  dit sasi, non bulan. La semaine des sept jours vient de l'arabe
%  par le malais, mais elle a pris en javanais une forme a elle, et
%  c'est cette forme-la qu'on ecrit.
%
%  ET LA SALUTATION (68) EST LE SEUL ENDROIT DU LIVRET OU LA REGLE DE
%  NIVEAU SE RETOURNE. La colonne tient le ngoko ; mais le petit
%  garcon poli salue des adultes qu'il connait, et le javanais N'A PAS
%  de salutation en ngoko a dire a un adulte. « sugeng enjing » est du
%  krama, et c'est la forme JUSTE : un enfant qui saluerait en ngoko
%  ne serait pas poli, il serait insolent, et l'alinea dit justement
%  qu'il est poli. Le niveau ne se choisit pas une fois pour toutes,
%  il se choisit par qui parle a qui. kolonoj.py ne releve ni sugeng
%  ni enjing, pour la meme raison qu'il ne releve pas dhahar.
%
%  LE CONTROLE A TROUVE SON PROPRE FAUX AMI DANS CE TABLEAU, ET IL A
%  FALLU LE DESARMER. Les oies (25) se disent « banyak » en javanais.
%  Or « banyak » etait releve comme mot-outil indonesien — il y
%  signifie « beaucoup », et le javanais dit « akeh ». Les deux mots
%  s'ecrivent pareil et ne sont pas le meme mot. La regle a ete otee
%  de kolonoj.py, avec sa raison ecrite a la place ou elle etait :
%  « banyak » javanais est un OISEAU, et un controle qui crie sur un
%  mot juste finit desarme. C'est la troisieme fois que cette lecon
%  sert dans le releve, et la premiere ou elle sert CONTRE une regle
%  qu'on venait d'ecrire soi-meme.
%
%  ET LE MEME PIEGE ETAIT ANNONCE AU TABLEAU 1 SOUS UN AUTRE MOT :
%  « tangga », le voisin en javanais et l'echelle en indonesien.
%  Deux faux amis en trois tableaux, tous les deux entre ces deux
%  langues-la : la ressemblance est justement ce qui rend le releve
%  difficile.
%
%  UNE FAUTE D'ORDRE, LA PREMIERE DE LA COLONNE, ET ELLE N'ETAIT PAS
%  STRUCTURELLE. Au bloc t03-c2-07-1 le feu d'artifice (20) avait ete
%  pose avant les roseaux (19), alors que l'ido dit « preparas inter
%  la kani la pirotekno ». Le javanais est a tete initiale et possede
%  « kang », relative postposee : rien dans la langue n'obligeait a
%  intervertir. C'etait un simple choix de tour, et c'est ce qui rend
%  la faute instructive — les colonnes a tete finale se trompent
%  parce que la grammaire les pousse, celle-ci s'est trompee sans
%  raison, et seule la machine pouvait le voir.
%
%  LE VETEMENT DE FETE EST NEERLANDAIS D'UN BOUT A L'AUTRE : klompen
%  les sabots (75), du neerlandais klomp ; dhasi la cravate (9), de
%  das ; krah le col (13), de kraag ; bretel les bretelles (15) ; pet
%  la casquette (11) ; rompi le gilet (17). Le portugais tient les
%  pieces plus anciennes — sepatu la chaussure (1, 5), de sapato ;
%  kemeja la chemise (12), de camisa ; gendera le drapeau (29), de
%  bandeira. Deux couches, deux siecles, et elles se partagent la
%  garde-robe par ordre d'arrivee.
%
%  LE PASAR MALEM (39) EST LE MOT JUSTE ET NON UNE TRANSPOSITION. La
%  fete foraine de la planche — chevaux de bois, cirque, tir, geant
%  et nain, menagerie — est exactement ce que le javanais appelle un
%  pasar malem, et le mot est celui de la chose, non celui d'une
%  chose d'ici qu'on lui substituerait. En revanche le bal public
%  (42) reste « pesta dansa » et non « jogedan » : le joged est une
%  danse javanaise, la planche grave un bal de sous-prefecture, et
%  c'est le refus deja pose au tableau 2 pour le gobak sodor.
% ===================================================================

%%K t03-apar-1 apar
\VUsaut{3.0mm}
\VUtitre{15.0pt}{60}{GAMBAR No. 3}
\VUsaut{1.6mm}
\VUtitre{11.4pt}{0}{Bocahan lan Nom-noman.}
\VUsaut{3.4mm}

%%K t03-apar-2 apar
(Gambar kaping 3 lan kaping 4 disusun mangkene supaya nyuguhake ing
patang \VUgras{adegan kulawarga} kawruh baku bab \VUgras{cuaca,}
\VUgras{umur,} \VUgras{kulawarga,} \VUgras{sandhangan} lan
\VUgras{kaanan.})

\VUsaut{4.0mm}
%%K t03-tit-1 sub
\VUcentre{12.0pt}{0}{\textit{Adegan Kapisan.}}
\VUsaut{3.0mm}

%%K t03-tit-2 sub
\VUpk{10.2pt}{40}{Upacara baptis ing desa.}
\VUsaut{2.4mm}

%%K t03-tit-3 sub
\VUpk{10.2pt}{40}{Mangsa semi.}
\VUsaut{3.0mm}

%%K t03-c1-01-1 p
1. --- Kita ana ing \VUgras{mangsa semi,} ing sasi \VUgras{Maret,}
\VUgras{April} utawa \VUgras{Mei,} ing dina-dinane \VUgras{minggu,}
\VUgras{Senen,} \VUgras{Selasa} utawa \VUgras{Rebo.} Wektu iku
\VUgras{jam sepuluh} esuk.

%%K t03-c1-02-1 p
2. --- \VUgras{Cuaca} apik, \VUgras{hawa} resik. Srengenge sumunar.
Sawetara \VUgras{mendhung cilik} \textsuperscript{(2)} munggah ing
\VUgras{langit} \textsuperscript{(1)} kang biru.

%%K t03-c1-03-1 p
3. --- \VUgras{Wit-witan} meh ora duwe \VUgras{godhong.}
\VUgras{Wit poplar} \textsuperscript{(3)} kang lencir lan \VUgras{wit
platan} \textsuperscript{(17)} kang dhuwur nganti saiki isih mung duwe
semi.

%%K t03-c1-04-1 p
4. --- \VUgras{Manuk sriti} \textsuperscript{(4)} padha bali lan mabur
mrana-mrene karo cuit-cuit bungah ing sakubenge \VUgras{pucuk
lancip} \textsuperscript{(7)} ing \VUgras{menara lonceng}
\textsuperscript{(5)} kang dadi makutha \VUgras{greja}
\textsuperscript{(6)} desa iku.

%%K t03-tit-4 sub
\VUsaut{3.4mm}
\VUpk{10.2pt}{40}{Greja.}
\VUsaut{3.0mm}

%%K t03-c1-05-1 p
5. --- Prasaja banget greja iku, karo \VUgras{gapura gedhene}
\textsuperscript{(13)} lan \VUgras{jam gedhene} \textsuperscript{(8)}.
\VUgras{Jarum cilik} \textsuperscript{(9)} ana ing angka X, kang
nuduhake \VUgras{jam.} \VUgras{Jarum gedhe} \textsuperscript{(10)} ana
ing XII — tengah wengi — kang nuduhake \VUgras{menit.}

%%K t03-c1-06-1 p
6. --- Ing menara, \VUgras{lonceng-lonceng} \textsuperscript{(11)}
muni kanthi bungah. Ing pucuking payon ngadeg \VUgras{salib}
\textsuperscript{(12)} watu cilik.

%%K t03-c1-07-1 p
7. --- Greja iku ora mung duwe \VUgras{gapura gedhe}
\textsuperscript{(13)}, nanging uga duwe lawang cilik ing sisih.
Metu saka lawang cilik iku \VUgras{romo pastur} \textsuperscript{(15)}
kang nganggo \VUgras{jubahe,} saka \VUgras{kamar sakristi}
\textsuperscript{(14)} tumuju \VUgras{pastoran} \textsuperscript{(19)}.

%%K t03-c1-08-1 p
8. --- Ing cedhake, iki \VUgras{bakul susu} \textsuperscript{(24)} karo
\VUgras{kuldine} \textsuperscript{(23)}. Dheweke bali saka kutha,
panggonan dheweke nggawa susu becik kanggo bocah-bocah.

%%K t03-tit-5 sub
\VUsaut{4.0mm}
\VUpk{10.2pt}{40}{Kebon woh-wohan.}
\VUsaut{3.4mm}

%%K t03-c1-09-1 p
9. --- Kang Aliboron — si kuldi — marem banget karo lakon esuk iku.
Dheweke ngambu kanthi nikmat hawa kang wangi dening ambune
\VUgras{kembang} \textsuperscript{(20)} kang nutupi \VUgras{wit
woh-wohan} \textsuperscript{(18)} ing \VUgras{kebon woh-wohan}
sandhinge. Abang mawar lan putih \VUgras{wit persik}
\textsuperscript{(18)} lan \VUgras{wit aprikot} \textsuperscript{(19)}
iku, lan padha njanjekake panen kang becik.

%%K t03-c1-10-1 p
10. --- Sauntara iku, \VUgras{tawon} \textsuperscript{(22)} cilik saka
\VUgras{omah tawon} \textsuperscript{(21)} padha mrana golek
\VUgras{madu} kang legi karo ngungingung.

%%K t03-c1-11-1 p
11. --- Nanging apa kang kedadeyan? Lah bocah-bocah desa padha
mlayu teka lan mbuyarake \VUgras{banyak} \textsuperscript{(25)} kang
kaget. Iku metune rombongan saka greja sawise \VUgras{baptis}e anake
notaris.

%%K t03-c1-12-1 p
12. --- Iakobus cilik, ing \VUgras{ayunan bayine}
\textsuperscript{(26)}, tangi amarga lonceng muni bungah, lan
mbakyune Magdalena, kang uga kepengin ndeleng arak-arakan baptis,
njaluk marang ibune supaya nyisiri rambute lan nyandhangi dheweke ing
undhak-undhakan omah.

%%K t03-c1-13-1 p
13. --- Ibune nyarujuki. Ibune nganggo \VUgras{kacu guluné}
\textsuperscript{(28)}, \VUgras{bluse} \textsuperscript{(29)} lan
\VUgras{clemeke} \textsuperscript{(30)}, banjur lungguh ing kursi cilik
ngarep lawang. Magdalena mung nganggo \VUgras{rok jerone}
\textsuperscript{(31)} lan \VUgras{korsete} \textsuperscript{(32)}.

%%K t03-c1-14-1 p
14. --- Bungah temen bocah wadon iku! Dheweke lali karo kembang
kesenengane, \VUgras{kembang anyelire} \textsuperscript{(34)} kang
diencoki \VUgras{kupu} \textsuperscript{(35)}, \VUgras{kembang
tulipe} \textsuperscript{(36)}, \VUgras{kembang Meine}
\textsuperscript{(37)} ing \VUgras{pot} \textsuperscript{(38)} ing
\VUgras{tembok} \textsuperscript{(33)}, lan \VUgras{kembang maware}
\textsuperscript{(39)} kang dadi pager endah temen. Dheweke namatake
sandhangan sarwa sugih para nyonyah ing rombongan.

%%K t03-c1-15-1 p
15. --- Bocah-bocah desa ora patiya nggatekake sandhangan. Padha
grubyug lan sesenggolan supaya oleh \VUgras{permen}
\textsuperscript{(55)} utawa \VUgras{dhuwit receh} \textsuperscript{(52)}.
Salah sijine \VUgras{nangis} \textsuperscript{(40)}.

%%K t03-c1-16-1 p
16. --- Kang liyane ngidak sikile \VUgras{asu} \textsuperscript{(42)}
kang mlayu karo nggonggong lan mbuyarake \VUgras{babon}
\textsuperscript{(43)} sakuthuk-\VUgras{kuthuke} \textsuperscript{(44)}.

%%K t03-tit-6 sub
\VUsaut{5.0mm}
\VUpk{10.2pt}{40}{Arak-arakan baptis.}
\VUsaut{4.0mm}

%%K t03-c1-17-1 p
17. --- Ing ngarepe rombongan mlaku \VUgras{ibu susu}
\textsuperscript{(45)}. Dheweke nganggo \VUgras{kupluk renda}
\textsuperscript{(46)} endah kanthi pita dawa. Dheweke nggendhong
\VUgras{bayi kang lagi lair} \textsuperscript{(47)} kang kabeh
sandhangane \VUgras{renda} \textsuperscript{(48)}.

%%K t03-c1-18-1 p
18. --- Ing mburine ibu susu teka \VUgras{bapak baptis}
\textsuperscript{(49)} lan \VUgras{ibu baptis} \textsuperscript{(53)}.
Wong loro iku mbagi permen lan dhuwit receh. Bapak baptis nganggo
\VUgras{sarung tangan} \textsuperscript{(51)} lan \VUgras{topi
silinder} \textsuperscript{(50)}. Ibu baptis nyekel \VUgras{karangan
kembang} \textsuperscript{(54)} endah ing tangane.

%%K t03-c1-19-1 p
19. --- Ing mburine, kita weruh \VUgras{pakdhe} \textsuperscript{(56)}
lan \VUgras{budhe} \textsuperscript{(58)} bayi iku. Budhene nganggo
\VUgras{topi} \textsuperscript{(59)} kang apik, pakdhene nganggo
\VUgras{jas dawa} \textsuperscript{(57)} ireng lan rompi putih. Iki
banjur \VUgras{sedulur lanang} \textsuperscript{(60)} lan
\VUgras{sedulur wadon} \textsuperscript{(61)}. Kang wadon nyekel
tangane \VUgras{adhine wadon} \textsuperscript{(62)} si bayi, yaiku
Berta cilik.

%%K t03-c1-20-1 p
20. --- \VUgras{Adhine lanang} \textsuperscript{(63)} Berta mlaku ing
mburi. Dheweke nyekel tangane \VUgras{sanak wadon}
\textsuperscript{(64)}. \VUgras{Kanca-kancane} \textsuperscript{(65)}
kulawarga teka ing mburine maneh.

%%K t03-tit-7 sub
\VUsaut{4.6mm}
\VUpk{10.2pt}{40}{Wong kang padha nonton.}
\VUsaut{3.4mm}

%%K t03-c1-21-1 p
21. --- Ing cedhake rombongan, iki \VUgras{wong ngemis}
\textsuperscript{(66)} kang nyanggong ing \VUgras{cagak keteke}
\textsuperscript{(67)}. Dheweke njaluk-njaluk.

%%K t03-c1-22-1 p
22. --- Ing sisih liya ana bocah lanang cilik kang \VUgras{sopan}
banget \textsuperscript{(68)}. Dheweke ngucap salam lan muni
« \VUgras{sugeng enjing} » marang wong-wong kang ditepungi.

%%K t03-c1-23-1 p
23. --- Wong desa padha metu saka omahe kanggo ndeleng arak-arakan
baptis. Kita weruh \VUgras{wong desa} \textsuperscript{(69)} karo
\VUgras{kupluk katune,} lan ing sandhinge \VUgras{wong wadon desa}
\textsuperscript{(70)}. Banjur \VUgras{wong wadon tuwa}
\textsuperscript{(71)} kang nyanggong ing \VUgras{tekene}
\textsuperscript{(72)}. Pungkasane \VUgras{tukang panggilingan}
\textsuperscript{(73)} karo \VUgras{topi wole} \textsuperscript{(74)}
kang amba, lan wong wadon desa enom kang nganggo \VUgras{klompen}
\textsuperscript{(75)}, kang lagi mulang anake mlaku.

%%K t03-c1-24-1 p
24. --- Adegan iku nyenengake banget.

\VUsaut{4.0mm}
%%K t03-tit-8 sub
\VUcentre{12.0pt}{0}{\textit{Adegan Kapindho.}}
\VUsaut{3.0mm}

%%K t03-tit-9 sub
\VUpk{10.2pt}{40}{Pesta umum.}
\VUsaut{2.4mm}

%%K t03-tit-10 sub
\VUpk{10.2pt}{40}{Dolanan banyu lan balapan prau.}
\VUsaut{3.0mm}

%%K t03-c2-01-1 p
1. --- \VUgras{Mangsa panas} wis teka. Srengenge kang panas banget ing
sasi \VUgras{Juni} (Juli utawa \VUgras{Agustus}) ngajak para nom-noman
adus lan ndayung.

%%K t03-c2-02-1 p
2. --- Ing pinggir tengene kali, para tukang dayung padha nyopot
sandhangan supaya bisa melu dolanan banyu lan balapan prau.

%%K t03-c2-03-1 p
3. --- Salah sijine nyopot \VUgras{sepatune} \textsuperscript{(1)} ;
dheweke nguculi taline. Kanca loro liyane wis siyap. Saiki
wong-wong iku mung nganggo \VUgras{klambi rajut} \textsuperscript{(2)}
lan \VUgras{kathok adus} \textsuperscript{(3)}. Padha ngobrol karo
ngenteni kanca-kancane. Padha ngomong bab preinan lan pesta kang
dianakake ing pinggir sabrang.

%%K t03-c2-04-1 p
4. --- Ing lemah, ing cedhake, gumlethak \VUgras{sandhangan}
kanca-kancane : sapasang \VUgras{sepatu bot} \textsuperscript{(4)},
\VUgras{sepatu} \textsuperscript{(5)}, \VUgras{kaos kaki}
\textsuperscript{(6)}, \VUgras{topi wol} \textsuperscript{(7)} lan
\VUgras{topi damen} \textsuperscript{(8)}, sarta \VUgras{dhasi}
\textsuperscript{(9)}.

%%K t03-c2-05-1 p
5. --- Salah sijine nom-noman nglempit \VUgras{setelane}
\textsuperscript{(10)} kanthi ngati-ati. Setelan iku diselehake ing
lawon, panggonan kancane sedhela maneh bakal ngempakake
\VUgras{sandhangan} loro-lorone.

%%K t03-c2-06-1 p
6. --- Bocah kaping nem kasep. Dheweke isih nganggo \VUgras{pet}
\textsuperscript{(11)} ing sirahe. Dheweke durung nyopot
\VUgras{kemejane} \textsuperscript{(12)}, \VUgras{krahe}
\textsuperscript{(13)}, utawa \VUgras{mansete} \textsuperscript{(14)}.
Dheweke nyekel \VUgras{rompine} \textsuperscript{(17)} ing tangan lan
isih nganggo \VUgras{kathok dawane} \textsuperscript{(16)} lan
\VUgras{bretele} \textsuperscript{(15)}. Mesthi dheweke ora bakal
siyap kanggo balapan kang bakal teka.

%%K t03-c2-07-1 p
7. --- Ing cedhake para nom-noman, ana dalan cilik kang ing
pinggir-pinggire akeh \VUgras{kembang eri} \textsuperscript{(18)},
tumuju menyang pinggir kali, panggonan Kang Iakobus nyiyapake
ing antarane \VUgras{glagah} \textsuperscript{(19)} \VUgras{mercon}
\textsuperscript{(20)} kang bakal diunekake mengko bengi.

%%K t03-tit-11 sub
\VUsaut{4.4mm}
\VUpk{10.2pt}{40}{Balapan.}
\VUsaut{3.4mm}

%%K t03-c2-08-1 p
8. --- Ing kali, sawetara nyonyah, kang ngeyubi awake nganggo payung,
mlaku-mlaku nganggo \VUgras{prau} \textsuperscript{(21)}. Padha ngetutake
balapan kang narik banget iku. Ing saben \VUgras{prau balap}
\textsuperscript{(22)}, \VUgras{tukang dayung} \textsuperscript{(24)}
papat ndayung kanthi sakuwate, dene \VUgras{juru mudhi}
\textsuperscript{(26)} nyekel \VUgras{kemudhi} \textsuperscript{(25)}
lan nuntun prau kang tipis iku. \VUgras{Dayung} \textsuperscript{(23)}
nggebug banyu kanthi ajeg lan prau kang entheng iku mlaku kanthi
cepet gawe mumet.

%%K t03-c2-09-1 p
9. --- Sawetara nom-noman rebutan, ing cedhake cagak kreteg, hadiah
ing \VUgras{cagak lunyu} \textsuperscript{(28)}. Salah sijine mlaku
kanthi ngati-ati ing cagak kang lentur lan lunyu iku supaya tekan
\VUgras{gendera} \textsuperscript{(29)} cilik kang dipasang ing
pucuke. \VUgras{Saingane} \textsuperscript{(30)} kang lagi apes tiba
ing banyu. Dheweke nglangi supaya tekan dharatan.

%%K t03-c2-10-1 p
10. --- Nom-noman kang luwih tuwa \VUgras{adu prau}
\textsuperscript{(32)} ing cedhake \VUgras{dhermaga}
\textsuperscript{(33)}. Kabeh dolanan iku ditonton \VUgras{wong akeh}
\textsuperscript{(31)} kang nyanggong ing pager tembok
\VUgras{kreteg} \textsuperscript{(27)} lan kang gembruduk ing pinggir
kali. Nanging sawetara wong kang nonton malah minger kanggo ngetutake
dolanan \VUgras{cagak hadiah} \textsuperscript{(45)}, utawa kanggo
nyawang \VUgras{rombongan wali kutha} kang teka kanggo \VUgras{mbagi}
hadiah.

%%K t03-c2-11-1 p
11. --- Kang kapisan, iki sawetara \VUgras{bocah}
\textsuperscript{(35)} kang mlaku ndhisiki \VUgras{para pemain musik}
\textsuperscript{(36)} ; banjur teka \VUgras{wali kutha}
\textsuperscript{(37)}, para pembantune lan \VUgras{dewan kutha}
kutha cilik iku. Pungkasane mlaku \VUgras{tukang pompa geni}
\textsuperscript{(38)}.

%%K t03-tit-12 sub
\VUsaut{5.0mm}
\VUpk{10.2pt}{40}{Pasar malem.}
\VUsaut{3.6mm}

%%K t03-c2-12-1 p
12. --- Wong akeh banget ana ing \VUgras{papan pasar malem}
\textsuperscript{(39)}. Wong padha teka ndeleng warna-warna
\VUgras{tarub} : \VUgras{jaran kayu} \textsuperscript{(40)},
\VUgras{sirkus} \textsuperscript{(41)} kang pagelarane wis diwiwiti ;
\VUgras{pesta dansa} \textsuperscript{(42)}, panggonan nom-noman lanang
lan wadon kutha ngrasakake senenge \VUgras{jejogedan.}

%%K t03-c2-13-1 p
13. --- Ing mburi, kita weruh \VUgras{arena} \textsuperscript{(44)},
panggonan \VUgras{matador} Spanyol kang kendel adu karo
\VUgras{banteng} \textsuperscript{(45)} kang ngamuk ; banjur
\VUgras{sepur luncur} \textsuperscript{(49)} lan \VUgras{papan
nembak} \textsuperscript{(46)}, panggonan wong latihan nganggo
\VUgras{bedhil} lan nembak \VUgras{sasaran.}

%%K t03-c2-14-1 p
14. --- Ing cedhake, iki \VUgras{panggung pasar malem}
\textsuperscript{(47)}, panggonan wong main \VUgras{komedi} lan
\VUgras{drama.} Cedhak banget ana tarube \VUgras{wong raksasa}
\textsuperscript{(48)} lan \VUgras{wong cebol.} \VUgras{Dhuwure} kang
kapisan rong meter luwih limalas senti ; kang kapindho meh padha
karo bocah cilik.

%%K t03-c2-15-1 p
15. --- Pungkasane, iki \VUgras{kandhang kewan} \textsuperscript{(50)}
gedhe karo \VUgras{kewan galake,} \VUgras{macan lorenge}
\textsuperscript{(51)} kang \VUgras{kukune} kuwat, \VUgras{singane}
\textsuperscript{(52)} kang \VUgras{surine} kandel, \VUgras{jerapahe}
\textsuperscript{(53)} loro lan \VUgras{gajahe} \textsuperscript{(54)}
kang trampil temen nganggo \VUgras{tlale.}

%%K t03-c2-16-1 p
16. --- \VUgras{Juru latih kewan} \textsuperscript{(55)} kang nglatih
kewan-kewan iku ana ing lawange tarub.

\VUsaut{6.0mm}
\VUfilet{22mm}
